\chapter{Introduction}
\label{chap:introduction}
% 1.1 in Kara & Garcia

The first quasar discovered, 3C 273 \parencite{Edge19593C273Catalogue}, had been documented in photographic plates at least as early as 1887 \parencite{Slavchevamihova20133c273half}. This emblematic radio-loud object showed variable emission in the optical and radio bands on timescales from days to years \parencites{SmithH19633C273Var2}{Dent19653C273Var}, indicating that the size of the emitting region spanned only a few light-days in size. Its compactness, coupled with the finding of its distant redshift \parencite{Schmidt19633C273z}, was very difficult to explain by dynamics that operate in normal galaxies, and variability thus became one of the key arguments that the enormous luminosities of quasars had to originate from an exceptionally compact and energetic central engine \parencites{Salpeter1964AccretionTheory}{Zeldovich1964AccretionTheory}.
% or some objects, rapid optical variability implied brightness temperatures that clearly required a nonthermal emission mechanism. For example, Oke (1967) observed day-to-day changes of 0.25 and 0.1 mag for 3C 279 and 3C 446 (Shields NED)

Now we know that variability in all wavelengths is an ubiquitous property of active galactic nuclei (AGN), the broader class of which quasars are the most luminous members. Its study has been essential in understanding the physics of their nuclear regions, which, in general, cannot be resolved by current telescopes. This had been strictly true only until the observations of Sgr A* and M87 \parencite{EventHorizonTelescope2019} and, more recently, with the GRAVITY interferometer on the VLT resolving the outer BLR in 3C 273 \parencite{Gravity20183C273BLR}; still, these surveys are limited to massive sources in the very nearby universe. 
For the vast majority of AGN, variability studies therefore rely on the \textit{light curve}, a time series recording how the brightness (or flux) of the AGN, measured within a given wavelength band, varies over a monitoring period. The time scales, spectral changes, and the correlations and delays between variations extracted from these light curves provide crucial information on the nature and location of these components, described in Chapter \ref{chap:agn}, and on their interdependencies \parencite{Ulrich1997VARIABILITY}. 

While variability-based methods are, generally speaking, more effective than other methods (e.g., SED-based) at providing a clearer signature of the nuclear activity against host galaxy emission \parencite{Yu2025ScalableEztao}, deriving the features and their uncertainties from an irregularly sampled light curve is very challenging; thereon, high quality light curves are indispensable to enable its robust characterization \parencite{Kelly2014CARMA}. 

The challenge is now pressing because of the scale of modern time-domain surveys. Wide-field surveys such as the Zwicky Transient Facility (\cite{Bellm2019ZTF}; \S\,\ref{ssec:ZTF}), and the forthcoming Legacy Survey of Space and Time (\cite{Ivezi2019LSST}; \S\,\ref{ssec:LSST}), monitor large areas of the sky repeatedly and deliver optical light curves for enormous numbers of AGN at once. These light curves are irregularly sampled and interrupted by gaps from the day-night cycle, weather, and the seasonal visibility of each field. Reconstructing the signal across these gaps is an underdetermined problem, since many different light curves can fit the same sparse observations; the problem therefore calls for a probabilistic method, one that returns a distribution of plausible light curves rather than a single estimate.

This work investigates the use of \textit{neural stochastic flows} \parencite{Kiyohara2025SNF}, a recently proposed class of generative models that learn the transition density of a stochastic process directly, as a flexible method for reconstructing AGN optical light curves. The aim is not to show that such a method surpasses existing approaches, but to establish whether it can be applied to this problem, to validate it against the analytic properties of the damped random walk commonly used to model the variability, and to characterize its behaviour as the simulated data are made progressively more realistic and it is finally applied to real light curves from the Zwicky Transient Facility.

This thesis is organized as follows. Chapter \ref{chap:agn} reviews the structure of active galactic nuclei and the emission from each of their components. Chapter \ref{chap:variability} describes the physical origin of AGN variability and the tools used to characterize it, introducing the damped random walk and the time-domain surveys relevant to this work. Chapter \ref{chap:simulation} sets out the simulation framework, deriving the transition probability of the damped random walk and defining the survey regimes reproduced in the simulations. Chapter \ref{chap:nsf} presents the reconstruction method, developing neural stochastic flows and the bridge construction used to fill the gaps. Chapter \ref{chap:reconstruction} reports the results on simulated and real ZTF light curves, and Chapter \ref{chap:conclusions} summarizes the findings and outlines future work.
